\documentclass[11pt,a4paper]{article}
\usepackage[margin=28mm]{geometry}
\usepackage{amsmath,amssymb,amsthm,booktabs}
\usepackage[T1]{fontenc}
\usepackage{lmodern}
\usepackage[colorlinks=true,urlcolor=blue,citecolor=blue,linkcolor=blue]{hyperref}
\newtheorem{theorem}{Theorem}
\newtheorem{lemma}[theorem]{Lemma}
\newtheorem{proposition}[theorem]{Proposition}
\newtheorem{remark}[theorem]{Remark}
\newcommand{\F}{\mathbb F_2}
\newcommand{\HH}{\mathcal H}
\newcommand{\verts}{\operatorname{V}}
\title{Explicit cocycles for graph Hall algebras}
\author{Lluis Eriksson}
\date{4 November 2026\\\small Research draft, candidate 2: external assessment pending}
\begin{document}
\maketitle
\begin{abstract}
For the reduced simplicial set whose simplices are subgraphs of a finite labelled
undirected multigraph equipped with ordered partitions of their vertices, we construct
explicit mod-two cohomology classes indexed by subgraphs. Vertex restriction
pulls a supported top-dimensional cochain into a global cocycle. Permutation
cycles pair with these classes as a Kronecker basis, while a vertex filtration
bounds the total cohomology dimension using augmented Boolean order complexes.
Alexander--Whitney multiplication recovers the graph Hall product. This proves
the finite-graph Hall/cohomology identification stated as
Conjecture 1.1 by Toury. The construction concerns labelled subgraphs, not
isomorphism classes of rooted forests. Exact finite checks accompany the proof;
priority and external mathematical assessment remain pending.
\end{abstract}

\section{The graph model and the statement}
Let $G=(V,E)$ be a finite labelled undirected graph. The finite edge set has
an endpoint map to unordered vertex pairs, allowing repeated endpoints (loops)
and distinct labelled edges with equal endpoint pairs. A subgraph $H$ specifies
$S\subseteq V$ and any subset of ambient edge identities with both endpoints
in $S$.
It need not be induced in $G$. Write $H|_A$ for the induced restriction of $H$
to $A\subseteq V$, and include the empty graph.

We use the simplicial graph model of Bergner et al.\ \cite{wald}, also described
in \cite{bergner,toury}. The ordered-partition description and its symmetry
also appear in \cite[Examples 3.9 and 5.11]{cms}; those constructions are
background, rather than claims of novelty here. In degree $k$, a simplex is a subgraph $H$ together with
an ordered partition $(B_1,\ldots,B_k)$ of $\verts(H)$, permitting empty blocks.
Degree zero has one element. Endpoint faces delete the first or last block,
restricting $H$ to the remaining vertices. Face $d_i$, $0<i<k$, merges $B_i$
and $B_{i+1}$. Degeneracies insert an empty block. Thus a normalized simplex
has all its blocks nonempty, and its degree is at most $|V|$.

The graph Hall algebra $\HH(G;\F)$ has a basis $[H]$ indexed by all actual
subgraphs. Its unit is $[\varnothing]$. If $A,B$ intersect in a vertex, set
$[A][B]=0$. Otherwise set
\begin{equation}\label{hallproduct}
 [A][B]=\sum_{D\subseteq E_G(\verts(A),\verts(B))}[A\amalg B+D].
\end{equation}
Here the summand retains exactly the internal edges of $A$ and $B$ and adds
$D$ as cross edges. This is the labelled graph Hall product
\cite[Proposition 4.4.1]{toury}; give $[H]$ degree $|\verts(H)|$.

\begin{theorem}\label{main}
There is a canonical graded algebra isomorphism
\[
 \Phi_G:\HH(G;\F)\longrightarrow H^*(|X_G|;\F).
\]
For a nonempty $H$ on $S$, its image is represented by the pullback, under
vertex restriction to $S$, of the cochain supported on one singleton-block
ordering of $H$. Its cohomology class is independent of that ordering.
\end{theorem}

This statement addresses the finite-graph statement of \cite[Conjecture 1.1]{toury}.
It identifies a cohomology algebra, not a homotopy type or the even-degree
polyhedral-product space constructed in that paper. No assertion is made about
arbitrary reduced 2-Segal sets. In particular, the construction does not pass
to isomorphism classes of unlabelled graphs or to a rooted-tree cut algebra.

\section{A finite dimension bound}
All chains and cochains below are normalized and have coefficients in $\F$.
Let $F_p C_*(X_G)$ be generated by simplices whose underlying graph has at
most $p$ vertices. This is a finite filtration by subcomplexes: interior faces
preserve the graph and endpoint faces decrease its vertex set.

\begin{lemma}\label{dimension}
The sum of the Betti numbers of $X_G$ is at most the number of subgraphs of $G$.
\end{lemma}
\begin{proof}
For $p>0$, the quotient $F_pC_*/F_{p-1}C_*$ is a direct sum over subgraphs
$H$ with exactly $p$ vertices. Each summand has the ordered partitions of
$\verts(H)$ into $k$ nonempty blocks as its degree-$k$ basis. Its differential
merges adjacent blocks; endpoint faces vanish in the quotient.

When $p\ge2$, associate $(B_1,\ldots,B_k)$ to the chain of proper nonempty
subsets
\[
 B_1\subsetneq B_1\cup B_2\subsetneq\cdots
 \subsetneq B_1\cup\cdots\cup B_{k-1}.
\]
This identifies the summand with the augmented simplicial chains of the
proper Boolean order complex, shifted up by two degrees. For $k=1$ the
chain of subsets is empty and represents the augmentation in degree $-1$.
Removing a proper initial union merges the corresponding adjacent blocks,
so the differentials agree over $\F$.

The proper Boolean order complex is the barycentric subdivision of the
boundary of the $(p-1)$-simplex. Its reduced homology is one-dimensional
in degree $p-2$ and zero otherwise. The shifted summand therefore has one
homology generator, in degree $p$. For $p=1$ this conclusion follows directly
from the single degree-one partition. For $p=0$, the empty graph contributes
the basepoint in degree zero.

For a short exact sequence of finite chain complexes, the long exact sequence
in homology bounds the total homology dimension of the middle complex by the
sum of those of the other two. Apply this successively to the vertex
filtration. Each subgraph contributes one to the resulting upper bound.
Normalized simplicial homology computes the homology of the realization;
over a field its Betti numbers equal those of cohomology.
\end{proof}

\section{Cycles and dual cocycles}
For a nonempty subgraph $H$ on $S$, put $p=|S|$. Define
\begin{equation}\label{cycle}
 z_H=\sum_{\sigma\in\operatorname{Ord}(S)}
       (H;\{\sigma_1\},\ldots,\{\sigma_p\}),
\end{equation}
where $\operatorname{Ord}(S)$ denotes all linear orderings of $S$.

\begin{lemma}\label{cycles}
The chain $z_H$ is a cycle.
\end{lemma}
\begin{proof}
Every internal face merges two adjacent singleton blocks. Exactly two
orderings yield that face, differing by swapping those vertices. Endpoint
faces also pair: for every $v\in S$ and every ordering of $S\setminus\{v\}$,
placing $v$ at the beginning or at the end yields the same induced graph
$H|_{S\setminus\{v\}}$ and the same residual ordering. All pairs cancel over
$\F$. For $p=1$, the two endpoints both equal the basepoint and cancel.
\end{proof}

For a fixed $S\subseteq V$, define vertex restriction on all simplices,
including degenerate ones, by
\[
 \rho_S(H;B_1,\ldots,B_k)
  =(H|_S;B_1\cap S,\ldots,B_k\cap S).
\]
Its target is $X_{G|_S}$. Intersection commutes with adjacent union, endpoint
deletion and insertion of empty blocks. Hence $\rho_S$ is a simplicial map.
On normalized cochains a restricted simplex with an empty block contributes zero.

Choose an ordering $\omega$ of $S$. In degree $p$ of $X_{G|_S}$, let
$a_{H,\omega}$ be the cochain taking value one precisely on
$(H;\{\omega_1\},\ldots,\{\omega_p\})$ and zero on every other normalized
$p$-simplex. As the target has no nondegenerate simplices in degree $p+1$,
this is a cocycle. Set
\begin{equation}\label{cocycle}
 \alpha_{H,\omega}=\rho_S^*a_{H,\omega}.
\end{equation}
Explicitly, it is one on a simplex precisely when its restriction to $S$
is that chosen ordered singleton simplex of $H$.

\begin{proposition}\label{basis}
For subgraphs $H,K$ with the same number of vertices,
\[
 \langle\alpha_{H,\omega},z_K\rangle=\begin{cases}1&H=K,\\0&H\ne K.\end{cases}
\]
The classes $[\alpha_{H,\omega}]$, together with the degree-zero unit, form
a basis of the cohomology of $X_G$.
\end{proposition}
\begin{proof}
A nonzero evaluation requires $S=\verts(H)$ to be contained in $\verts(K)$.
Equal cardinality forces equality of these vertex sets. It then requires
exactly the same edge set and the chosen ordering. Precisely one term of
\eqref{cycle} contributes when $H=K$. Thus the displayed pairing is a
Kronecker pairing. Since a coboundary vanishes on every cycle, the classes
are independent. Their total number is the upper bound in
Lemma~\ref{dimension}, so they are a basis. For $G$ empty, both sides have
only the degree-zero unit.
\end{proof}

\begin{lemma}\label{order}
The class $[\alpha_{H,\omega}]$ is independent of $\omega$.
\end{lemma}
\begin{proof}
It suffices to exchange an adjacent pair in an ordering of $S$. For $p\ge2$,
in $X_{G|_S}$ take the degree-$(p-1)$ cochain supported on $H$ with that
pair merged into a doubleton block and every other block unchanged. Its
coboundary is the sum of the two ordered top cochains. Indeed, a normalized
$p$-simplex there must have all $p$ vertices in singleton blocks; only the
two splittings of the specified doubleton yield the supported face.
Endpoint faces have fewer vertices and do not contribute. Pulling back
proves the claim. When $p=1$ there is only one ordering.
\end{proof}

\section{Multiplication}
Write $\alpha_H$ for a representative with an arbitrary chosen order,
and use brackets when passing to its order-independent class.

\begin{proof}[Proof of Theorem~\ref{main}]
Define $\Phi_G([H])=[\alpha_H]$ for nonempty $H$ and send the empty graph
to the degree-zero unit. Proposition~\ref{basis} makes this a graded vector
space isomorphism. It remains to check multiplication.

Let $A,B$ have vertex sets $S,T$ and degrees $p,q>0$. The
Alexander--Whitney product evaluates $\alpha_A$ on the front $p$-face and
$\alpha_B$ on the back $q$-face of a $(p+q)$-simplex. These two faces
retain respectively its first $p$ and last $q$ vertex blocks. If $S\cap T$
is nonempty, a shared vertex would have to occur in both disjoint groups
of blocks. Thus the cochain product is identically zero.

Suppose $S\cap T=\varnothing$. Restrict the simplex to $S\cup T$. For a
nonzero product its first $p$ restricted blocks must be the ordered
singletons of $A$, and its last $q$ restricted blocks the ordered
singletons of $B$. The graph restricted to $S$ must equal $A$, and that
restricted to $T$ must equal $B$. The only unconstrained edges of the
restricted graph are ambient edges between $S$ and $T$. Consequently there
is a cochain identity
\[
 \alpha_A\smile\alpha_B
 =\sum_{D\subseteq E_G(S,T)}\alpha_{A\amalg B+D,\,\omega_S\omega_T},
\]
where the order on the right is concatenation and each cochain is pulled
back from $X_{G|_{S\cup T}}$. This identity also accounts for additional
ambient vertices: they disappear under restriction and cannot change
the required singleton blocks. Lemma~\ref{order} allows replacement of
concatenated orders by the chosen orders in cohomology. This gives exactly
\eqref{hallproduct}. The degree-zero unit acts as the constant cochain one.
Therefore $\Phi_G$ is an algebra isomorphism. The absence of a choice in
its cohomology classes proves the asserted canonicity.
\end{proof}

\section{Finite verification and scope of evidence}
The companion files \texttt{graph\_chains.py} and \texttt{verify.py} run with
Python 3.10 or later, without third-party dependencies, using
\texttt{python verify.py}. The exact code hashes and captured output are in
the accompanying SHA-256 manifest and \texttt{verification-results.json}.
The additional script \texttt{verify\_extension.py} and
\texttt{extension-results.json} check loops and repeated endpoint pairs.
The three-vertex mixed fixture has one loop at vertex 0, two labelled edges
between 0 and 1, and two between 1 and 2.
The supplied programs use normalized chain matrices and exact elimination over
$\F$. They check $d^2=0$, cocycle closure, cycle pairings and
Alexander--Whitney multiplication modulo coboundaries. They supply controls
for the construction, not a replacement for its proof.
\begin{center}
\begin{tabular}{lrr}
\toprule
Graph & Subgraph cocycles & Cup products checked\\
\midrule
$P_1$ & 1 & 0\\
$P_2$ & 4 & 4\\
$P_3$ & 12 & 39\\
$P_4$ & 33 & 265\\
Triangle & 17 & 45\\
Four-vertex star & 34 & 273\\
Four-cycle & 46 & 324\\
$K_4$ & 112 & 512\\
Four isolated vertices & 15 & 132\\
One loop & 2 & 0\\
Two labelled loops & 4 & 0\\
Two parallel edges & 6 & 4\\
Two loops, two parallel edges & 20 & 16\\
Three-vertex mixed graph & 50 & 128\\
\bottomrule
\end{tabular}
\end{center}
The product counts include positive degrees whose sum is at most the vertex
count. Larger degrees vanish by dimension. The unit is handled in the proof,
not in those pair loops. The code and proof were developed in the same
assistant context; the computations do not constitute independent review.

The proof applies to finite labelled undirected multigraphs with loops over $\F$.
For a fixed subgraph $H$, the filtration summand depends only on its vertex set.
Restriction retains precisely the labelled edge identities whose endpoints
survive. Thus the cycle cancellations, cocycles and order exchange do not
require simplicity. In the product argument every retained edge is internal
to one of the two disjoint vertex sets or crosses between them. Loops are
internal and parallel edge identities remain separate choices. This partitions
the allowed edge sets exactly as in \eqref{hallproduct}.
We assert no integral or odd-characteristic version, identification of
homotopy types, or result about interacting quantum fields. The theorem
must be distinguished from Toury's already constructed even-degree
polyhedral-product realization. The exact graph model and literal
labelled-subgraph Hall multiplication are essential.

\section{Model assessments and review context}
This draft was developed with AI assistance in a Codex chat, including
literature triage, proof construction and verification code. The same
assistant performed the internal checks. No separate model assessment of
this exact PDF, human peer review, formal verification, editorial acceptance
or importance rating has yet been obtained. No stars are assigned here.
Requested model names or a product label do not authenticate the effective
model identity or reasoning effort. Targeted bibliographic searches have
not established worldwide priority. Lluis Eriksson is also AIRR's founder, registry operator and responsible
editor (\url{https://airr.science/governance/}); this is a disclosed
author--operator conflict. No editorial decision is recorded in this draft.
Earlier small graph computations and
the formulations of the problem are credited to the sources below.

\begin{thebibliography}{9}
\bibitem{cms} I. Contreras, R. A. Mehta and W. H. Stern,
\emph{Frobenius and commutative pseudomonoids in the bicategory of spans},
arXiv:2311.15342v2 (2024), Examples 3.9 and 5.11.
\url{https://arxiv.org/abs/2311.15342}.
\bibitem{wald} J. E. Bergner, A. M. Osorno, V. Ozornova, M. Rovelli and
C. I. Scheimbauer, \emph{2-Segal sets and the Waldhausen construction},
Topology and its Applications 235 (2018), 445--484.
\href{https://doi.org/10.1016/j.topol.2017.12.009}{doi:10.1016/j.topol.2017.12.009}.
\bibitem{bergner} J. E. Bergner, \emph{Combinatorial examples and applications
of 2-Segal sets}, arXiv:2411.18544v1 (2024).
\url{https://arxiv.org/abs/2411.18544}.
\bibitem{toury} L. Toury, \emph{Hall algebras of graphs and rooted trees},
arXiv:2609.16849v1 (2026), Conjecture 1.1 and Sections 4.1, 4.7.
\url{https://arxiv.org/abs/2609.16849}.
\bibitem{owr} J. Bergner (joint work with N. Kuhn and I. Zakharevich),
\emph{2-Segal sets, algebraic K-theory, and Hall algebras}, in
Homotopy Theory, Oberwolfach Report 36/2019.
\url{https://doi.org/10.4171/OWR/2019/36}.
\end{thebibliography}
\end{document}
